% vorspann.tex – Kompetenzblatt, zusammenbau v0.7 (2026-09-28).
% Ergänzt mathblatt.sty für Kompetenzblätter, ohne sie zu ändern
% (layout-befunde.md 1–34 vom 28.09.). Wird nach \usepackage{mathblatt}
% eingelesen.
\makeatletter
\RequirePackage{newunicodechar}
% Zeichen, die Latin Modern nicht hat (Befund 20): Text aus einer
% Ersatzschrift, Mathe als Befehl; ohne Ersatzschrift auch Text als Mathe.
\newif\ifkbersatz
\IfFontExistsTF{FreeSerif}{\newfontfamily\kbersatzschrift{FreeSerif}\kbersatztrue}{%
 \IfFontExistsTF{DejaVu Serif}{\newfontfamily\kbersatzschrift{DejaVu Serif}\kbersatztrue}{%
  \IfFontExistsTF{Cambria Math}{\newfontfamily\kbersatzschrift{Cambria Math}\kbersatztrue}{%
   \IfFontExistsTF{Segoe UI Symbol}{\newfontfamily\kbersatzschrift{Segoe UI Symbol}\kbersatztrue}{}}}}
\newcommand{\kbzeichen}[2]{\ifmmode #2\else\ifkbersatz{\kbersatzschrift #1}\else\ensuremath{#2}\fi\fi}
\newcommand{\kbzeichenhoch}[1]{\ifkbersatz\ifmmode\text{\kbersatzschrift #1}\else{\kbersatzschrift #1}\fi\else ?\fi}
% Kopf- und Fußzeile (Befund 1, 2): Kopf Thema · Blattart, Fuß Kennung und Seite
\newcommand{\kbkopfzeile}[2]{\pagestyle{fancy}\fancyhf{}\renewcommand{\headrulewidth}{0pt}%
  \fancyhead[L]{\footnotesize\color{mbgrau}#1}%
  \fancyfoot[L]{\footnotesize\color{mbgrau}#2}%
  \fancyfoot[R]{\footnotesize\color{mbgrau}Seite \thepage}}
% Flattersatz überall (Befund 25)
\raggedright
% Titel des Blatts: Ich-kann-Satz, darunter Zweigzeile
\newcommand{\kbtitel}[2]{\par\noindent{\Large\bfseries\raggedright #1\par}%
  \ifblank{#2}{}{\par\smallskip\noindent{\small #2}\par}\medskip}
% Abschnitt: „Das kennst du schon“, „Zum Merken“
\newcommand{\kbabschnitt}[1]{\par\addvspace{12pt}\Needspace*{5\baselineskip}%
  \noindent{\bfseries #1}\par\nobreak\vspace{3pt}}
% Hauptnummer: Titel hängt am ersten Block; jeder Block (Teilaufgabe) bricht
% nicht, passt er nicht mehr, rückt er ganz auf die nächste Seite (Befund 21).
\newif\ifkbtitel
\newsavebox{\kbbox}
\newlength{\kbhoehe}
\newlength{\kbnummerabstand}\setlength{\kbnummerabstand}{10pt}
\newlength{\kbteilabstand}\setlength{\kbteilabstand}{6pt}
\newenvironment{kbaufgabe}[1]{\stepcounter{aufgabe}\setcounter{teil}{0}%
  \gdef\kbtiteltext{\noindent\hangindent1.6em\hangafter1\makebox[1.6em][l]{\textbf{\theaufgabe.}}#1\par}%
  \global\kbtiteltrue\par\addvspace{\kbnummerabstand}}{\par}
\newenvironment{kbblock}{\par\addvspace{\kbteilabstand}\begin{lrbox}{\kbbox}%
  \begin{minipage}[t]{\linewidth}\raggedright\parskip\z@
  \ifkbtitel\kbtiteltext\vspace{2pt}\global\kbtitelfalse\fi
  \list{}{\leftmargin1.6em\labelwidth1.4em\labelsep0.2em\topsep\z@\partopsep\z@
    \parsep\z@\itemsep\z@\rightmargin\z@}\raggedright}%
  {\endlist\end{minipage}\end{lrbox}%
  \setlength{\kbhoehe}{\dimexpr\ht\kbbox+\dp\kbbox\relax}%
  \Needspace*{\kbhoehe}\noindent\usebox{\kbbox}\par}
% Paarsatz (Platz nutzen, Befund 27): zwei Hauptnummern oder zwei
% Teilaufgaben mit Zeichenaufgabe nebeneinander, als ein Block
\newcommand{\kbnummer}[1]{\stepcounter{aufgabe}\setcounter{teil}{0}%
  \noindent\hangindent1.6em\hangafter1\makebox[1.6em][l]{\textbf{\theaufgabe.}}#1\par\vspace{2pt}}
\newcommand{\kbhalb}[1]{\list{}{\leftmargin1.6em\labelwidth1.4em\labelsep0.2em\topsep\z@
  \partopsep\z@\parsep\z@\itemsep\z@}\raggedright #1\endlist}
\newcommand{\kbpaar}[2]{\par\addvspace{\kbnummerabstand}\begin{lrbox}{\kbbox}%
  \begin{minipage}[t]{0.485\linewidth}\raggedright\parskip\z@ #1\end{minipage}\hfill
  \begin{minipage}[t]{0.485\linewidth}\raggedright\parskip\z@ #2\end{minipage}\end{lrbox}%
  \setlength{\kbhoehe}{\dimexpr\ht\kbbox+\dp\kbbox\relax}\Needspace*{\kbhoehe}\noindent\usebox{\kbbox}\par}
\newcommand{\kbteilpaar}[2]{\item[]\noindent
  \begin{minipage}[t]{0.485\linewidth}\kbhalb{#1}\end{minipage}\hfill
  \begin{minipage}[t]{0.485\linewidth}\kbhalb{#2}\end{minipage}\par}
% Anweisung der Hauptnummer, eigene Zeile über der Teilaufgabe (Befund 5)
\newcommand{\kbanweisung}[1]{\item[]#1\par\vspace{2pt}}
% Kopfzeile einer Teilaufgabe: Buchstabe, Auftakt (halbfett oder Kontext),
% Prüfkennung klein rechts (Befund 3, 12, Beschluss c)
\newlength{\kbkennbreite}\setlength{\kbkennbreite}{1.9cm}
\newcommand{\kbkennung}[1]{{\footnotesize #1}}
\newcommand{\kbteil}[3]{\item[#1]\ifblank{#3}%
  {\parbox[t]{\linewidth}{\raggedright #2}}%
  {\parbox[t]{\dimexpr\linewidth-\kbkennbreite\relax}{\raggedright #2}\hfill
   \parbox[t]{\kbkennbreite}{\raggedleft\kbkennung{#3}}}\par}
% Frage in eigener Zeile (Befund 4), ein Satz je Zeile (Befund 10)
\newcommand{\kbfrage}[1]{\par\vspace{2pt}#1\par}
% Antwortfeld in eigener Zeile darunter, linksbündig (Befund 4, 8)
\newcommand{\kbantwort}[1]{\par\vspace{4pt}\noindent#1\par}
% Zwei Spalten: links Grafik/Tabelle/Aufgabe, rechts Antwort oder Rechenraum
% (Befund 27); oben bündig. Beginnt einen eigenen Listenpunkt ohne Marke.
\newcommand{\kbzweispaltig}[3][0.48]{\item[]\noindent
  \begin{minipage}[t]{#1\linewidth}\raggedright #2\end{minipage}\hfill
  \begin{minipage}[t]{\dimexpr0.98\linewidth-#1\linewidth\relax}\raggedright #3\end{minipage}\par}
% Linke Spalte mit Teilaufgabe (Buchstabe hängt links heraus wie in der Liste)
\newcommand{\kbspalte}[1]{\list{}{\leftmargin\z@\labelwidth1.4em\labelsep0.2em\topsep\z@
  \partopsep\z@\parsep\z@\itemsep\z@}\raggedright #1\endlist}
% Antwortfelder: 2,2 cm, in Punkten 1,2 cm; ohne Umbruch davor
\newcommand{\kbfeld}[1][]{\mbox{\underline{\hspace{2.2cm}}\ifblank{#1}{}{\,#1}}}
\newcommand{\kbfeldk}[1][]{\mbox{\underline{\hspace{1.2cm}}\ifblank{#1}{}{\,#1}}}
% Grafik oder Tabelle unter dem Text, linksbündig (Befund 8, 28)
\newcommand{\kbdarunter}[1]{\par\vspace{4pt}\noindent#1\par}
% Ankreuzen: je Option eine Zeile, linksbündig (Befund 6, 7)
\newcommand{\kbkreuz}[1]{\par\vspace{2pt}\noindent$\square$\ #1\par}
% Bearbeitungsraum (Befund 9): Schreibzeilen wie mathblatt (9 mm)
\newcommand{\kbraum}[1]{\schreibzeilen{#1}}
% Wertetabelle mit Köpfen im Textmodus (Befund 26): wie \wertetabelle der
% Vorlage, aber die Köpfe setzt der Aufruf selbst ($x$, „Zeit in h“).
\NewDocumentCommand{\kbwertetabelle}{o m m m}{%
  \def\mbkopf{}\def\mbwerte{}\def\mbleer{}\setcounter{mbn}{0}%
  \foreach \v in {#4}{\xappto\mbkopf{& $\v$}\xappto\mbleer{& }\stepcounter{mbn}\xdef\mbnn{\arabic{mbn}}}%
  \IfNoValueTF{#1}{}{\foreach \v in {#1}{\ifdefempty{\v}{\xappto\mbwerte{& }}{\xappto\mbwerte{& $\v$}}}}%
  \begingroup\renewcommand{\arraystretch}{1.3}%
  \edef\mbpre{|>{\noexpand\columncolor{mbkasten}}l|*{\mbnn}{>{\noexpand\centering\noexpand\arraybackslash}p{\noexpand\mbzell}|}}%
  \expandafter\mbtabstart\expandafter{\mbpre}\hline
  #2 \mbkopf \\ \hline
  \IfNoValueTF{#1}{\rule{0pt}{\mbzeile}#3 \mbleer \\ \hline}{#3 \mbwerte \\ \hline}
  \end{tabular}\endgroup}
% Merkkasten am Ende (Aufbau des Kompetenzblatts)
\newcommand{\kbmerk}[2]{\par\addvspace{14pt}\begin{lrbox}{\kbbox}\begin{minipage}[t]{\linewidth}%
  \noindent{\bfseries #1}\par\vspace{4pt}%
  \uebersichtskasten{\raggedright #2}\end{minipage}\end{lrbox}%
  \setlength{\kbhoehe}{\dimexpr\ht\kbbox+\dp\kbbox\relax}\Needspace*{\kbhoehe}\noindent\usebox{\kbbox}\par}
% Lösungsblatt: Nummer und Buchstabe halbfett, Lösung daneben
\newcommand{\kbloesung}[2]{\par\addvspace{3pt}\noindent\hangindent2.8em\hangafter1\makebox[2.8em][l]{\textbf{#1}}#2\par}
% Lernblatt (zusammenbau v0.9): Antwortfeld hinter kurzen Termen,
% Verweis der Zone auf die Nummer im Lernblatt
\newcommand{\lbfeld}{\mbox{\underline{\hspace{2.6cm}}}}
% „Das kann ich“ unter „Prüfe dich“ (wie abhakseite, ohne neue Seite)
\newenvironment{lbdaskannich}{\par\addvspace{14pt}\Needspace*{8\baselineskip}%
  \einheitenkopf[abhaken][Das kann ich]{Das kann ich}%
  \begingroup\small\setlength{\parskip}{1pt}\setlength{\parindent}{0pt}}%
  {\par\endgroup}
\newcommand{\lbhaengst}[1]{\par\vspace{3pt}{\footnotesize\color{mbgrau}Hängst du hier $\rightarrow$ Nr.~#1}\par}
% Zwei Teilaufgaben nebeneinander; Buchstaben in derselben Flucht wie
% untereinander (links hängend)
\newlength{\lbhalbbreite}
\newcommand{\lbpaar}[2]{\item[]\hspace*{-\leftmargin}%
  \setlength{\lbhalbbreite}{\dimexpr0.5\linewidth+0.5\leftmargin-0.6em\relax}%
  \begin{minipage}[t]{\lbhalbbreite}\kbhalb{#1}\end{minipage}\hfill
  \begin{minipage}[t]{\lbhalbbreite}\kbhalb{#2}\end{minipage}\par}
% Lernblatt v1.0 (Befund 6 und 13 des Lehrers vom 01.10.): Striche für den
% Rechenweg so stark wie die Antwortlinie, halbe Breite, eine Zeile Luft
% zwischen Aufgabentext und erstem Strich bzw. Antwortfeld
\newlength{\lbstrichbreite}
\renewcommand{\kbraum}[1]{\par\setlength{\lbstrichbreite}{0.5\textwidth}%
  \ifdim\lbstrichbreite>\linewidth\setlength{\lbstrichbreite}{\linewidth}\fi
  \vspace{\baselineskip}\setcounter{mbsz}{0}%
  \loop\ifnum\value{mbsz}<#1\relax
    \nointerlineskip\vskip\mbschreibhoehe
    \hbox{\hskip\@totalleftmargin\underline{\hspace{\lbstrichbreite}}}%
    \stepcounter{mbsz}\repeat
  \nointerlineskip\hbox{\vrule\@width\z@\@height\z@\@depth\mbschreibtiefe}\prevdepth\z@}
\renewcommand{\kbantwort}[1]{\par\vspace{\baselineskip}\noindent#1\par}
% Antwort rechts neben der Grafik, etwa auf halber Höhe der Grafik
\newcommand{\lbantwortrechts}[2][0pt]{\par\vspace*{#1}\noindent#2\par}
% Kennzeichen „GYM“ rechts im Titel der Hauptnummer (wie die Prüfkennung)
\newcommand{\lbgym}{\hfill{\footnotesize GYM}}
% Lernblatt v1.1 (Befunde des Lehrers am TER-L5)
\RequirePackage{multicol}
\tcbuselibrary{breakable}
% Hauptnummer ohne Titel (Blatt 0): die Nummer steht links vor dem Buchstaben
\newenvironment{lbaufgabe}{\stepcounter{aufgabe}\setcounter{teil}{0}%
  \global\kbtitelfalse\par\addvspace{\kbnummerabstand}}{\par}
\newcommand{\lbnr}{\llap{\textbf{\theaufgabe.}\hspace{0.9em}}}
% Test „Kannst du das schon?“: hellgrauer Kasten ohne Nummer, Lösung klein
% unten rechts
\newcommand{\lbtest}[3]{\par\addvspace{8pt}\noindent
  \fcolorbox{gray!55}{gray!10}{\parbox{\dimexpr\linewidth-2\fboxsep-2\fboxrule\relax}{%
  \raggedright{\bfseries #1}\par\vspace{3pt}#2\par\vspace{2pt}%
  {\raggedleft\footnotesize Lösung: #3\par}}}\par\addvspace{6pt}}
\newcommand{\lbtt}[2]{\par\noindent\hangindent1.6em\hangafter1%
  \makebox[1.6em][l]{#1}#2\par\vspace{3pt}}
% „Verstanden?“ im Kasten wie der Test (bricht über die Seite)
\newenvironment{lbkasten}{\par\addvspace{8pt}\begin{tcolorbox}[breakable,
  colback=gray!10,colframe=gray!55,boxrule=0.4pt,arc=0pt,left=3pt,right=3pt,
  top=2pt,bottom=4pt]}{\end{tcolorbox}}
\newcommand{\lbhinweis}[1]{{\small\itshape #1}\par\vspace{2pt}}
% Lösungen: je Teilaufgabe eine Zeile, Nummer nur an der ersten
\newcommand{\lbloes}[3]{\par\noindent\hangindent2.8em\hangafter1%
  \makebox[1.5em][l]{\textbf{#1}}\makebox[1.3em][l]{#2}#3\par}
% Kopfzeile der Lösungsseiten fest „… · Lösungen“ (multicols überdeckt die
% Marke des Einheitenkopfs)
\newcommand{\lbkopfloesungen}{\renewcommand{\mbkopfeinheit}{\,\textperiodcentered\,Lösungen}}
\newcommand{\lbloesgrafik}[1]{\par\noindent\hspace*{2.8em}\resizebox{!}{2.2cm}{#1}\par}
\makeatother
